\documentclass[10pt]{article}

\usepackage[utf8]{inputenc}

\usepackage[T1]{fontenc}

\usepackage{amsmath}

\usepackage{amsfonts}

\usepackage{amssymb}

\usepackage[version=4]{mhchem}

\usepackage{stmaryrd}



\begin{document}

может быть приложена и к всевозможным С. с. п. в узком смысле вида



\[

Y_{\Phi}(s)=\Phi[X(t+s)]

\]



где $\Phi[X(t)]$ - произвольный функционал от С. с. п. $X(t)$, являющийся случайной величиной, имеюпцей математич. ожидание; если для всех таких С. с. п. $Y_{\Phi}(s)$ соответствующий предел $\hat{Y}_{\Phi}$ совпадает с E $Y_{\Phi}$, то С. с. п. $X(t)$ наз. ме т р и ч е с к и т р а в т и в н ы м. Для гауссовских С. с. п. $X(t)$, условие стационарности к-рых в узком смысле совпадаөт с условием стационарности в широком смысле, метрич. транзитивность будет иметь место тогда у только тогда, когда спектральная функция $F(\lambda)$ процесса $X(t)$ является непрерывной функцией $\lambda$ (см., нащр., [2], [3]). В общем случае нет простых необходимых б достаточных условий метрич. транзитивности С. с. п. $X(t)$.



Кроме указанного выше результата, касающегося. метрич. транзитивности, имеется также множество других результатов, относящихся к гауссовским С. с. П.; в частности, для таких процессов подробно изучен вопрос о локальных свойствах реализаций (т. е. отдельных наблюденных значений) С. с. ■. $X(t)$, о статистич. свойствах последовательности нулей или максимумов реализации С. с. п. $X(t)$ и точек пересечения ею заданного уровня (см., напр., [3]). Тұпичным примером результатов, касающихся пересечений уровня, является утверждение о том, что при широких условиях регулярности совокупность точек пересечения высокого уровня $x=u$ гауссовским С. с. Ш. $X(t)$ в нек-ром специальном масштабе времени (зависящем от $u$ и быстро стремящемся к бесконечности при $u \rightarrow \infty$ ) сходится при $u \rightarrow \infty$ к пуассоновскому потоку событийі единичной интенсивности (см [3]).



При рассмотрении С. с. ш. в широком смысле вводят в рассмотрение гильбертово пространство $H_{X}$ всевозможных линейных комбинаций значений процесса $X(t)$ и пределов в среднем квадратичном последовательностей таких линейных комбинаций, скалярное произведение в к-ром задается формулой $\left(Y_{1}, Y_{2}\right)=\mathrm{E} Y_{\mathbf{1}} \bar{Y}_{\mathbf{2}}$. В таком случае преобразование $X(t) \rightarrow X(t+a)$, где $a$ - фиксированное число, будет порождать линейный унитарный оператор $U_{a}$, отображающй̈ пространство $H_{X}$ на себя; при этом семейство операторов $U_{a}$ очевидно удовлетворяет условию $U_{a} U_{b}=U_{a+b}$, а значения $X(t)=U_{t} X(0)$ представляют собой совокупность точек (кривую, если время $t$ непрерывно, и счетную последовательность точек, если время дискретно), переводимую в себя всеми операторами $U_{a}$. Соответственно этому теория С. с. п. в широком смысле может быть переформулирована в терминах функционального анализа как изучение совокупностей точек $X(t)=U_{t} X(0)$ гильбертова пространства $H_{X}$, где $U_{t}$ семейство линейных унитарных операторов таких, что $U_{a} U_{b}=U_{a+b}$.



Центральное место в теории С. с. п. в широком смысле занимают спектральные рассмотрения, опирающиеся на разложение случайного процесса $X(t)$ и его корреляционной функции $B(\tau)$ в интеграл Фурье - Стилтьеса. В силу теоремы Хинчина [4] (являющейся простым следствием аналитич. теоремы Бохнера об общем виде положительно определенной функции) корреляционная функция $B(\tau)$ C. с. п. с непрерывным временем всегда может быть представлена в виде



\[

B(\tau)=\int_{\Lambda} e^{i \tau \lambda} d F(\lambda),

\]



где $F(\lambda)$ - ограниченная монотонно неубывающая функция $\lambda$, а $\Lambda=(-\infty, \infty)$; теорема Герглотца об общем виде положительно определенных последовательностей аналогичным образом показывае', тто такое же представление, но только с $\Lambda=[-\pi, \pi]$, имеет место и для корреляционной функции С. с. П. с дискретным временем. Если корреляционная функция $B(\tau)$ достаточно быстро убывает при $|\tau| \rightarrow \infty$ (как әто чаще всего и бывает в приложениях при условии, что под $X(t)$ понимается разность $X(t)-m$, т. е. считается, что $E X(t)=0)$, то интеграл в правой части (3) обращается в обыкновенный интеграл Фурье



\[

B(\tau)=\int_{\Lambda} e^{i \tau \lambda} \cdot f(\lambda) d \lambda

\]



где $f(\lambda)=F^{\prime}(\lambda)$ - неотрицательная функция. Функция $F(\lambda)$ наз. с пе к т р а льн о ї ф у нк п и е й C. с. п. $X(t)$, а функция $f(\lambda)$ (в случаях, когда имеет место равенство (4j) - его с п е к т р а л ь н о й п ло тн о с т ь ю. Исходя из формулы Хинчина (3) (или из задания С. с. п. $X(t)$ в виде совокупности точек $X(t)=$ $=U_{t} X(0)$ гильбертова пространства $H_{X}$ и теоремы Стоуна о спектральном представлении однопараметрич. групп унитарных операторов в гильбертовом пространстве), можно также показать, что сам процесс $X(t)$ допускает спектральное разложение вида



\[

X(t)=\int_{\Lambda} e^{i t \lambda} d Z(\lambda)

\]



где $Z(\lambda)-$ с лу ча й ная фу нкция с неко рре ли и о ванными пр и ра шения ми (т. е. такая, что $E d Z\left(\lambda_{1} \overline{d Z\left(\lambda_{2}\right)}=0\right.$ при $\left.\lambda_{1} \neq \lambda_{2}\right)$, удовлетворяющая условию $E|d Z(\lambda)|^{2}=d F^{\prime}(\lambda)$, а интеграл справа понимается как предел в среднем квадратичном соответствующей последовательности интегральных сумм. Разложение (5) дает основание рассматривать любой С. с. п. в широком смысле $X(t)$ как суперпозицию совокупности некоррелированных друг с другом гармонич. колебаний различных частот со случайным амплитудами и фазами; при этом спектральная функция $F(\lambda)$ и спектральная плотность $f(\lambda)$ определяют распределение средней энергии (или, точнее, мощности) входящих в состав $X(t)$ гармонич. колебаний по спект ру частот $\lambda$ (в связи с чем в прикладных исследованиях функция $f(\lambda)$ часто наз. такяе ә н е р г е т и ч ес к и с пек т о м, илис пе к т ром мо що ос т и, С с. П. $X(t))$.



Спектральное разложение корреляционної функции $B(\tau)$, задаваемое формулоӥ (3), показывает, что отображение $X(t) \rightarrow e^{i t \lambda,}$ переводящее элементы $X(t)$ гильбертова пространства $H_{X}$ в әлементы $e^{i t \lambda}$ гильбертова пространства $L^{2}(d F)$ комплекснозначных функций на множестве $\Lambda$ с интегрируемым по $d F(\lambda)$ квадратом модуля, является изометрич. отображением $H$ на $L^{2}(d F)$. Это отображение может быть далее продолжено до изометрического Јинеӥного отображения, $M$ всего пространства $H_{X}$ в пространство $L^{2}(d F)$, что позволяет переформулировать многие задачи теории С. с. п. в широком смысле в виде задач теории функций.



Значительная qасть теории С. с. п. в широком смысле посвящена методам решения линейных ашпроксимационных задач для таких процессов, т. е. методам нахождения линейной комбинации каких-то «известных» значений $X(t)$, к-рая лучше всего (в смысле минимума среднеквадратичной ошибки) приблиякает нек-рое «неизвестное» значение того же процесса или какую-то «неизвестную» случайную величину $Y$. В частности, задача об оптимальной линейной әкстраполяции C. с. II. $X(t)$ состоит в отыскании наилучшего приближения $X^{*}(s)$ к значению $X(s), s>0$, линейно зависящего от «прошлых значений» $X(t)$ с $t \leqslant 0$; задача об оптимальной линейной интерполяции - в отыскании наилучшего приближения $\kappa \quad X(s)$, линейно зависящего от значений $X(t)$, где $t$ пробегает все значения, не принадлежащие выделенному интервалу оси вре-





\end{document}